% =============================================================================
%  C2: Hierarchical viability composition — skeleton paper
%  STATUS: UNPROVEN. Known partial counterexample for RAFs.
% =============================================================================
\documentclass[11pt]{article}

\usepackage[utf8]{inputenc}
\usepackage[T1]{fontenc}
\usepackage[expansion=false]{microtype}
\usepackage{amsmath, amssymb, amsthm, mathtools}
\usepackage[margin=1in]{geometry}
\usepackage{enumitem}
\usepackage[hidelinks]{hyperref}

\theoremstyle{plain}
\newtheorem{theorem}{Theorem}
\newtheorem*{namedconjecture}{Conjecture}
\theoremstyle{definition}
\newtheorem{definition}[theorem]{Definition}

\title{Hierarchical Viability Composition \\
       \large (UNPROVEN --- skeleton paper for C2)}
\author{Liam McGuire\thanks{Unaffiliated researcher, Seattle, USA. liammcg44@gmail.com}}
\date{\today}

\begin{document}
\maketitle

\begin{abstract}
This paper is a \textbf{skeleton} that records C2 from the scaffolding
paper~\cite{emergent_systems_2026}. \textbf{No proof is attempted.}
A known partial counterexample (unions of RAFs are not always RAFs)
means the conjecture as stated needs to be refined into a precise
composition rule before any proof can land. The result is UNPROVEN.
\end{abstract}

\section{Statement (UNPROVEN)}

\begin{quote}
\emph{``C2 (hierarchical viability composition): under the entity
composition rule (E3), each viability formalism in C1 lifts to
hierarchical entities such that $F^*$-validity at the composite level
implies $F$-validity of each component.''}
\end{quote}
\noindent (Scaffolding paper~\cite{emergent_systems_2026},
\S\texttt{sec:c2}.)

\section{Plan of attack (no proof attempted)}

\begin{enumerate}[leftmargin=2em,itemsep=4pt]
  \item \textbf{Resolve the RAF counterexample (UNPROVEN).} Hordijk
    \& Steel~\cite{hordijk2004raf} showed unions of two RAFs are not
    in general RAFs. The refined composition rule must require joint
    catalyses or restrict to subsets where the union is preserved.
  \item \textbf{Autopoietic case (UNPROVEN).} Beer's gliders compose
    under simple unions; verifying that the composite-level closure
    operator is the join of component-level closures is the work.
  \item \textbf{Markov-blanket case (UNPROVEN).} Palacios et
    al.~\cite{palacios2020markov} proved Markov-blanket structure is
    closed under hierarchical nesting. This is essentially C2 for
    the MB case; citing and translating to the framework's vocabulary
    is the work.
  \item \textbf{MCC case (UNPROVEN; depends on C1).} If C1's
    reformulation of MCC as a closure operator goes through, C2 for
    MCC follows by a similar union-or-restriction argument.
  \item \textbf{Categorical statement (UNPROVEN).} The cleanest
    formulation uses the categorical closure operators of
    Dikranjan and Tholen; the conjecture then reads ``viability is
    a lax morphism of operad-algebras.''
\end{enumerate}

\section{Connection to LE}

C2 characterises which compositions of level-$k$ entities are valid
level-$(k+1)$ entities. It is the post-transition state of LE's
mechanism (\texttt{le\_level\_emergence/}); the two should be coordinated.

\section{Status and follow-up}

\textbf{Current status: UNPROVEN.} Depends on C1; coordinate with LE.

\bibliographystyle{plain}
\bibliography{references}

\end{document}
